\section{Experimental Design}
\label{sec:design}

\subsection{Paired Trials}

Every protocol is evaluated over 30 trials in each channel configuration, for
600 runs in the headline study. The trials are paired rather than independent.
Trial index $i$ seeds a single random stream that generates, in fixed order, the
node positions, the initial energies, the full per-link shadowing matrix and the
sensed-value array for all rounds. That order depends only on $i$ and never on
which protocol is being run, so at a given index all ten configurations face a
byte-identical physical world. A protocol cannot see a different topology by
virtue of being a different protocol.

Protocol-internal randomness, meaning LEACH's election draws, the genetic
operators and the $\varepsilon$-greedy exploration during training, comes from a
separate stream derived by a keyed hash of the run index and protocol name.
Python's built-in hash is salted per process and is never used anywhere in the
simulator, since it would silently destroy reproducibility across invocations.

The pairing is what licenses the statistical procedure in
Section~\ref{sec:stats}, and it is a substantially stronger design than the
unpaired mean-and-standard-deviation comparison this literature normally
reports. Topology variance is large relative to the differences between
protocols, and an unpaired test spends most of its power fighting variance that a
paired design eliminates by construction.

One thing is deliberately not shared: the per-packet channel draws. Different
protocols make different numbers of draws per round, so there is no coherent
sense in which two protocols could experience the same coin flips packet by
packet. The fairness-critical quantity is the static per-link shadowing, which is
identical. These are distinct claims, that the physical channel geometry is fair
and that the exact sequence of random outcomes is fair, and only the first needs
to hold.

\subsection{The Channel Sensitivity Pair}
\label{sec:channel-pair}

The entire study is run twice, once with the packet-error model of
Section~\ref{sec:channel} enabled and once with it disabled. This is not a
robustness appendix but a decision rule fixed before the results were read:

\begin{quote}
A conclusion that holds in both channel configurations is a statement about the
protocol. A conclusion that changes between them is a statement about transmit
power.
\end{quote}

Transmit power is a parameter the reference specifications do not fix, and
Section~\ref{sec:channel} showed that the packet-error curve is a waterfall
whose position determines how severely single-hop protocols are penalized on the
sink hop. Choosing that parameter to make a metric discriminating would be
precisely the calibration this study forbids. Running both endpoints and
reporting only what survives removes the choice from the argument.
Section~\ref{sec:results} reports one conclusion that the rule retracts.

\subsection{The Node-Count and Field-Area Sweep}
\label{sec:scale-design}

A single operating point cannot tell a reader whether a result is about
clustering or about one deployment. We therefore repeat the lossy study over a
$3\times3$ grid of node counts $N \in \{50, 100, 150\}$ and square field sides
$W \in \{50, 100, 150\}$~m, at 15 paired runs per cell, for a further 1{,}350
runs. Four design decisions are worth stating because each one bounds what the
sweep can conclude.

The base station scales with the field at $(W/2, 1.5W)$, which keeps the
deployment geometrically self-similar. It also means field size sets channel
severity as well as area: at $50\times50$~m the mean node-to-sink distance is
52~m and no node lies beyond 120~m, at $100\times100$~m the mean is 104~m with
33\% beyond, and at $150\times150$~m the mean is 156~m with 75\% beyond. That is
a real confound between ``bigger area'' and ``worse channel''. We report the mean
sink distance per cell so a reader can separate the two rather than engineering
the confound away and losing self-similarity instead.

The density radius feature scales with the field at $0.25W$. It is a
configuration field read identically by every protocol, so scaling it is a shared
environment choice rather than per-protocol tuning. Held fixed at 25~m it would
cover 78\% of a $50\times50$ field and 9\% of a $150\times150$ one, silently
changing what the feature means across the sweep.

The DQN and the GCN use their frozen $N = 100$, $100\times100$ weights in every
cell. No retraining happens anywhere in the sweep. This makes the sweep a
measurement of out-of-distribution transfer rather than of retrained
performance, and it is the reason the untrained fuzzy system matters as a
control: a protocol with no training at all cannot suffer transfer failure, so if
it moves the same way the learned ones do, the cause is the environment.

Everything else is unchanged from the headline study. Run index $i$ seeds the
topology, shadowing and sensed stream identically for every protocol within a
cell, so comparisons inside a cell stay paired. Comparisons across cells are not
paired, because a different $N$ means a different number of nodes to seed, and
they are not tested as if they were.

\subsection{Metrics}
\label{sec:metrics}

Lifetime is reported at three points, first, half and last node death, plus the
area under the living-node curve, which is the integral of surviving nodes over
rounds. All four are reported together because the three death points do not
agree on an ordering, for reasons Section~\ref{sec:results} makes concrete. Runs
that reach the round cap without full depletion are flagged as censored, and a
censored last-death figure is a lower bound rather than a measurement; it is
never averaged into a mean without the censoring count reported beside it.

Delivery is reported in two units, and the distinction matters more than it
appears to. Packets received at the sink measures aggregation depth, not
delivery: one chain round delivers a single fused packet where one clustered
round delivers roughly five, so on a packet basis the chain protocol appears five
times worse while actually delivering more sensor readings. Original readings
represented in delivered packets is the cross-protocol comparable quantity, and
the packet count is reported only as an aggregation-ratio indicator.

The same defect once affected the delivery ratio, which is why it is now
reading-based. The packet-based formula reported ten percent on a lossless
channel, purely because ten source readings became one aggregated packet.
Suppressed readings are not counted as losses in either numerator or
denominator, since a reading that was never injected was never lost. The cost of
suppression appears instead in data yield, defined as readings delivered over
readings taken, which is the metric that legitimately penalizes the reactive
protocols.

Computational cost is reported in three separate forms that are never merged:
measured wall-clock setup time and peak memory, counted algorithmic operations,
and analytical complexity. Wall-clock reflects implementation quality as much as
algorithm design; counted operations are more honest, and analytical complexity
is the most honest of the three.

\subsection{A Measurement Artifact Worth Reporting}

Timing and memory are measured in separate runs, and the reason generalizes
beyond this study. Allocation tracing intercepts every allocation and its
overhead grows with the number of traced blocks, so measuring both in one run
inflated setup time by factors between 3.5 and 7.8 depending on how many small
objects a protocol allocated. The self-organizing map, which allocates roughly
forty thousand temporary arrays per retraining, was hit hardest; the vectorized
heuristics barely at all.

Because the inflation was non-uniform, it distorted the ranking rather than only
the scale. The GCN's true setup cost is about 35\% higher than LEACH's, but under
tracing it measured about 30\% lower, so the sign of the comparison reversed. The
resolution is that trial zero of each protocol measures memory and reports no
timing, while the remaining trials do the reverse. Peak memory therefore comes
from a single sample and has no defined deviation, and timing means are over
$n-1$ trials. Simulation outcomes are unaffected, since tracing changes
wall-clock only and never energy or delivery. We report this because any
comparative study that profiles memory and time simultaneously is silently
ranking its subjects by allocation count.

\subsection{Statistical Procedure}
\label{sec:stats}

Because trials are paired, each comparison against a reference protocol reduces
to a vector of 30 per-topology differences. The null hypothesis that the
difference distribution is symmetric about zero is tested by a sign-flip
permutation test~\cite{good2005permutation}: 20{,}000 random sign vectors are
applied to the observed differences, and the $p$-value is the proportion of
permuted means at least as extreme as the observed one, with the standard
add-one correction to both numerator and denominator so that no $p$-value is ever
reported as exactly zero. The floor at 30 paired samples is therefore
$5\times10^{-5}$, and comparisons at that floor are reported as such rather than
as zero.

The procedure was validated against exhaustive enumeration on a reduced sample
where all $2^n$ sign assignments are enumerable, giving an exact $p$ of 0.00391
against a sampled 0.00420. Confidence intervals come from 20{,}000 paired
bootstrap resamples of the same difference vector~\cite{efron1993bootstrap}.
Because nine comparisons are made against the reference within each metric and
channel, all $p$-values are corrected by the Holm-Bonferroni step-down
procedure~\cite{holm1979stepdown}, which controls the family-wise error rate
without assuming independence between comparisons. The corrected floor is
therefore $p_{\mathrm{Holm}} = 0.00045$. Every statistic uses a fixed seed and is
reproducible exactly.

Permutation was chosen over a paired $t$-test because it assumes only
exchangeability under the null, not normality, and several of these difference
distributions are visibly skewed; the GCN's first-death figures are bimodal
across topologies. It was chosen over Wilcoxon signed-rank because it operates on
the actual differences rather than their ranks, so the effect size and the test
statistic are the same quantity~\cite{demsar2006statistical}.

\subsection{Verification}

Six gates were required to pass before any number in this paper was read.
Energy conservation is asserted at the end of every single run to $10^{-9}$~J
against the identity in~\eqref{eq:conservation}. Radio energies are checked
against hand-computed values on both sides of the crossover distance. A suite of
invariants confirms that energy never goes negative, that dead nodes never
transmit, that the living count is monotone, that the three death points are
ordered, and that delivery ratios remain in range. Determinism is verified by
running every protocol three times per seed and diffing all per-round and
summary fields, excluding only the timing and memory fields that are not
expected to be reproducible; all nine protocols reproduce exactly. A static
audit of the protocol source confirms that no implementation writes to the
energy or liveness arrays, calls the charging function, or chooses its own
packet sizes, across eighteen checks, all passing. Finally the aggregation
pipeline is verified by recomputing one protocol's summary row independently and
diffing at $10^{-9}$.

A seventh gate exists specifically because of the lossy channel. The
packet-error curve is emitted as a published artifact over 0 to 200~m and
asserted to be non-decreasing, effectively zero at intra-cluster range, and
strictly between zero and one somewhere in the band of sink distances. Had that
last condition failed, the channel would be behaving as a binary range switch
and the lossy results would carry no information.

The engine, both learned models' forward and backward passes, and every
statistic are implemented in NumPy~\cite{harris2020numpy}, with figures produced by
Matplotlib~\cite{hunter2007matplotlib} and no other numerical dependency.
The headline sweep of 600 runs completes in 80 minutes on sixteen cores, and the
1{,}350-run scale sweep in 222 minutes on fourteen. All configuration constants
live in a single frozen structure, no constant is defined anywhere else, and the
raw per-round output of every run is retained rather than only the summaries. The
central cell of the scale sweep reproduces the first 15 runs of the headline
lossy study exactly for all ten protocols, on both first node death and area
under the living-node curve, which is the check that the sweep harness and the
headline harness are the same simulator.
